\documentclass[11pt,a4paper]{article}

% Packages
\usepackage[utf8]{inputenc}
\usepackage[T1]{fontenc}
\usepackage{lmodern}
\usepackage{geometry}
\usepackage{graphicx}
\usepackage{listings}
\usepackage{xcolor}
\usepackage{hyperref}
\usepackage{booktabs}
\usepackage{fancyhdr}
\usepackage{tcolorbox}
\usepackage{enumitem}

% Page geometry
\geometry{margin=1in}

% Colors
\definecolor{codegreen}{rgb}{0,0.6,0}
\definecolor{codegray}{rgb}{0.5,0.5,0.5}
\definecolor{codepurple}{rgb}{0.58,0,0.82}
\definecolor{backcolour}{rgb}{0.95,0.95,0.92}
\definecolor{errorred}{rgb}{0.8,0,0}
\definecolor{warnyellow}{rgb}{0.7,0.5,0}
\definecolor{infogreen}{rgb}{0,0.5,0}
\definecolor{debugcyan}{rgb}{0,0.5,0.7}
\definecolor{tracemagenta}{rgb}{0.5,0,0.5}

% Code listing style
\lstdefinestyle{cstyle}{
    backgroundcolor=\color{backcolour},
    commentstyle=\color{codegreen},
    keywordstyle=\color{codepurple},
    numberstyle=\tiny\color{codegray},
    stringstyle=\color{codepurple},
    basicstyle=\ttfamily\small,
    breakatwhitespace=false,
    breaklines=true,
    captionpos=b,
    keepspaces=true,
    numbers=left,
    numbersep=5pt,
    showspaces=false,
    showstringspaces=false,
    showtabs=false,
    tabsize=4,
    language=C,
    morekeywords={size_t, uint8_t, int32_t, uint32_t, FILE, User, DEBUG_LEVEL,
                  DEBUG_TIMESTAMPS, DEBUG_COLORS, DEBUG_THREAD_SAFE, NDEBUG}
}

\lstset{style=cstyle}

% Header/footer
\pagestyle{fancy}
\fancyhf{}
\fancyhead[L]{\texttt{debug.h} API Reference}
\fancyhead[R]{Version 2.0}
\fancyfoot[C]{\thepage}

% Hyperref setup
\hypersetup{
    colorlinks=true,
    linkcolor=blue,
    filecolor=magenta,
    urlcolor=cyan,
}

% Title
\title{%
    \texttt{debug.h} \\
    \large A Battle-Hardened Debugging Header for C23/C++23 \\
    \vspace{0.5cm}
    \large API Reference \& User Guide
}
\author{GitHub Copilot \\ EdgeOfAssembly/RetroCodeMess}
\date{December 2025 \\ Version 2.0}

\begin{document}

\maketitle
\tableofcontents
\newpage

% ============================================================================
\section{Introduction}
% ============================================================================

\texttt{debug.h} is a feature-complete, header-only debugging library for C23 and C++23 that provides powerful debugging capabilities with zero runtime cost in release builds.

\subsection{Key Features}

\begin{itemize}
    \item \textbf{Five Log Levels}: TRACE, DEBUG, INFO, WARN, ERROR
    \item \textbf{Function Tracing}: Entry/exit logging with return values
    \item \textbf{Type-Safe Inspection}: Uses C11 \texttt{\_Generic} for automatic formatting
    \item \textbf{Cross-Platform}: Works on Linux, macOS, and Windows
    \item \textbf{Thread-Safe}: Optional mutex locking
    \item \textbf{Timestamps}: Optional date/time prefixes
    \item \textbf{ANSI Colors}: Configurable color output
    \item \textbf{Output Redirection}: File output and custom handlers
    \item \textbf{Zero Overhead}: Compiles out completely with \texttt{-DNDEBUG}
\end{itemize}

\subsection{Quick Start}

\begin{lstlisting}
#define DEBUG_LEVEL 5       // Full verbosity
#define DEBUG_TIMESTAMPS    // Add timestamps
#define DEBUG_COLORS        // Color output
#include "debug.h"

int main(void) {
    DEBUG_INFO("Hello, debugging world!");
    
    int x = 42;
    DEBUG_VAR(x);  // Output: VAR x = 42
    
    return 0;
}
\end{lstlisting}

% ============================================================================
\section{Configuration}
% ============================================================================

\subsection{Configuration Macros}

Define these \textbf{before} including \texttt{debug.h}:

\begin{table}[h]
\centering
\begin{tabular}{@{}lll@{}}
\toprule
\textbf{Macro} & \textbf{Default} & \textbf{Description} \\
\midrule
\texttt{DEBUG\_LEVEL} & 5 & Verbosity level (0--5) \\
\texttt{DEBUG\_TIMESTAMPS} & undefined & Enable timestamp prefixes \\
\texttt{DEBUG\_COLORS} & undefined & Enable ANSI color output \\
\texttt{DEBUG\_THREAD\_SAFE} & undefined & Enable mutex locking \\
\texttt{NDEBUG} & undefined & Disable all debug code \\
\bottomrule
\end{tabular}
\caption{Configuration macros}
\end{table}

\subsection{Log Levels}

\begin{table}[h]
\centering
\begin{tabular}{@{}clll@{}}
\toprule
\textbf{Level} & \textbf{Constant} & \textbf{Description} & \textbf{Color} \\
\midrule
1 & \texttt{LOG\_ERROR} & Errors only & \textcolor{errorred}{Red (bold)} \\
2 & \texttt{LOG\_WARN} & Warnings and above & \textcolor{warnyellow}{Yellow} \\
3 & \texttt{LOG\_INFO} & Info and above & \textcolor{infogreen}{Green} \\
4 & \texttt{LOG\_DEBUG} & Debug and above & \textcolor{debugcyan}{Cyan} \\
5 & \texttt{LOG\_TRACE} & Everything (verbose) & \textcolor{tracemagenta}{Magenta} \\
\bottomrule
\end{tabular}
\caption{Log levels and colors}
\end{table}

% ============================================================================
\section{Logging Macros}
% ============================================================================

\subsection{Basic Logging}

\begin{tcolorbox}[title=Logging Macros]
\begin{lstlisting}[numbers=none]
DEBUG_TRACE(fmt, ...);  // Most verbose
DEBUG_DEBUG(fmt, ...);
DEBUG_INFO(fmt, ...);
DEBUG_WARN(fmt, ...);
DEBUG_ERROR(fmt, ...);  // Least verbose
DEBUG_LOG(level, fmt, ...);  // Explicit level
\end{lstlisting}
\end{tcolorbox}

\textbf{Example:}
\begin{lstlisting}
DEBUG_INFO("Processing item %d of %d", current, total);
DEBUG_WARN("Cache miss for key: %s", key);
DEBUG_ERROR("Failed to open file: %s (errno=%d)", filename, errno);
\end{lstlisting}

\textbf{Output:}
\begin{verbatim}
[INFO ] main.c:10 process_items: Processing item 5 of 100
[WARN ] cache.c:42 lookup: Cache miss for key: user_123
[ERROR] file.c:15 open_file: Failed to open file: data.txt (errno=2)
\end{verbatim}

% ============================================================================
\section{Function Tracing}
% ============================================================================

\subsection{Entry and Exit}

\begin{tcolorbox}[title=Function Tracing Macros]
\begin{lstlisting}[numbers=none]
DEBUG_ENTER(fmt, ...);  // Log function entry
DEBUG_EXIT(ret);        // Log exit and return value
\end{lstlisting}
\end{tcolorbox}

\textbf{Example:}
\begin{lstlisting}
int factorial(int n) {
    DEBUG_ENTER("(n=%d)", n);
    DEBUG_ASSERT(n >= 0, "Negative input!");
    
    if (n <= 1) {
        DEBUG_EXIT(1);
    }
    
    int result = n * factorial(n - 1);
    DEBUG_VAR(result);
    DEBUG_EXIT(result);
}
\end{lstlisting}

\textbf{Output:}
\begin{verbatim}
[TRACE] math.c:2 factorial: ENTER (n=5)
[TRACE] math.c:2 factorial: ENTER (n=4)
[TRACE] math.c:2 factorial: ENTER (n=3)
[TRACE] math.c:2 factorial: ENTER (n=2)
[TRACE] math.c:2 factorial: ENTER (n=1)
[TRACE] math.c:6 factorial: EXIT (return: 1)
[TRACE] math.c:10 factorial: VAR result = 2
[TRACE] math.c:11 factorial: EXIT (return: 2)
...
\end{verbatim}

% ============================================================================
\section{Variable Inspection}
% ============================================================================

\subsection{DEBUG\_VAR}

The \texttt{DEBUG\_VAR} macro uses C11 \texttt{\_Generic} to automatically format variables based on their type.

\begin{tcolorbox}[title=Supported Types]
\begin{itemize}[noitemsep]
    \item \texttt{int}, \texttt{unsigned int}
    \item \texttt{long}, \texttt{unsigned long}
    \item \texttt{double}
    \item \texttt{char*}
    \item \texttt{void*}
\end{itemize}
\end{tcolorbox}

\textbf{Example:}
\begin{lstlisting}
int count = 42;
double pi = 3.14159;
char *name = "Alice";
void *ptr = malloc(100);

DEBUG_VAR(count);  // VAR count = 42
DEBUG_VAR(pi);     // VAR pi = 3.141590
DEBUG_VAR(name);   // VAR name = Alice
DEBUG_VAR(ptr);    // VAR ptr = 0x7f...
\end{lstlisting}

\subsection{DEBUG\_PTR}

For pointers with size descriptions:

\begin{lstlisting}
char *buffer = malloc(1024);
DEBUG_PTR(buffer, "1024 bytes");
// Output: PTR buffer @ 0x7f... (size: 1024 bytes)
\end{lstlisting}

% ============================================================================
\section{Memory Allocation}
% ============================================================================

\begin{tcolorbox}[title=Memory Macros]
\begin{lstlisting}[numbers=none]
void* ptr = DEBUG_MALLOC(size);
DEBUG_FREE(ptr);
\end{lstlisting}
\end{tcolorbox}

\textbf{Example:}
\begin{lstlisting}
typedef struct {
    int id;
    char name[32];
} User;

User *user = (User*)DEBUG_MALLOC(sizeof(User));
// Output: MALLOC 0x55... (size: 40 bytes)

// ... use user ...

DEBUG_FREE(user);
// Output: FREE 0x55...
\end{lstlisting}

% ============================================================================
\section{Assertions}
% ============================================================================

\subsection{DEBUG\_ASSERT}

Hard assertion that aborts on failure with backtrace:

\begin{lstlisting}
void *ptr = malloc(1024);
DEBUG_ASSERT(ptr != NULL, "Allocation failed for %zu bytes", 1024);
\end{lstlisting}

\subsection{DEBUG\_CHECK}

Soft check that logs warning but continues:

\begin{lstlisting}
DEBUG_CHECK(value >= 0, "Warning: negative value %d", value);
// Program continues even if check fails
\end{lstlisting}

% ============================================================================
\section{Backtrace}
% ============================================================================

\begin{lstlisting}
DEBUG_BACKTRACE();
\end{lstlisting}

\textbf{Output (Linux):}
\begin{verbatim}
[ERROR] main.c:42 risky_function: BACKTRACE (5 frames):
[ERROR]   #0: ./prog(risky_function+0x37) [0x55...]
[ERROR]   #1: ./prog(main+0x15c) [0x55...]
[ERROR]   #2: /lib/x86_64-linux-gnu/libc.so.6(+0x2a1ca) [0x7f...]
[ERROR]   #3: /lib/x86_64-linux-gnu/libc.so.6(__libc_start_main+0x8b)
[ERROR]   #4: ./prog(_start+0x25) [0x55...]
\end{verbatim}

\textbf{Note:} Compile with \texttt{-rdynamic} for function names.

% ============================================================================
\section{Loop Tracing}
% ============================================================================

\begin{lstlisting}
DEBUG_LOOP_START("Processing %d items", count);
for (int i = 0; i < count; i++) {
    DEBUG_LOOP_ITER(i, count);
    process(items[i]);
}
DEBUG_LOOP_END("Completed");
\end{lstlisting}

\textbf{Output:}
\begin{verbatim}
[TRACE] process.c:10 process_items: LOOP START Processing 100 items
[DEBUG] process.c:12 process_items: LOOP ITER 0/100
[DEBUG] process.c:12 process_items: LOOP ITER 1/100
...
[TRACE] process.c:15 process_items: LOOP END Completed
\end{verbatim}

% ============================================================================
\section{Output Redirection}
% ============================================================================

\subsection{File Output}

\begin{lstlisting}
FILE *logfile = fopen("debug.log", "w");
DEBUG_SET_OUTPUT(logfile);

DEBUG_INFO("This goes to the file");
// ... more logging ...

DEBUG_RESET_OUTPUT();  // Back to stderr
fclose(logfile);
\end{lstlisting}

\subsection{Custom Handler}

\begin{lstlisting}
void my_handler(const char* level, const char* file,
                int line, const char* func, const char* msg) {
    // Custom logging logic (e.g., send to network, database)
    syslog(LOG_INFO, "[%s] %s:%d %s: %s", level, file, line, func, msg);
}

DEBUG_SET_HANDLER(my_handler);
\end{lstlisting}

% ============================================================================
\section{Thread Safety}
% ============================================================================

Enable thread-safe logging:

\begin{lstlisting}
#define DEBUG_THREAD_SAFE
#include "debug.h"
\end{lstlisting}

Compile with: \texttt{-pthread}

\textbf{Note:} Uses \texttt{pthread\_mutex\_t} in C, \texttt{std::mutex} in C++.

% ============================================================================
\section{GDB Integration}
% ============================================================================

\subsection{DEBUG\_BREAK}

Insert a breakpoint:

\begin{lstlisting}
if (suspicious_condition) {
    DEBUG_BREAK();  // Triggers SIGTRAP
}
\end{lstlisting}

\subsection{Debugging Tips}

\begin{enumerate}
    \item Compile with \texttt{-g} for symbols
    \item Use \texttt{-rdynamic} for backtrace function names
    \item In GDB: \texttt{catch signal SIGTRAP} to catch \texttt{DEBUG\_BREAK}
    \item Use \texttt{bt} in GDB to see call stack
\end{enumerate}

% ============================================================================
\section{Compilation}
% ============================================================================

\subsection{Debug Build}

\begin{verbatim}
gcc -std=c23 -g -O0 -Wall -Wextra -rdynamic -o prog prog.c
\end{verbatim}

\subsection{Release Build}

\begin{verbatim}
gcc -std=c23 -O2 -DNDEBUG -o prog prog.c
\end{verbatim}

\subsection{Thread-Safe Build}

\begin{verbatim}
gcc -std=c23 -g -O0 -Wall -Wextra -rdynamic -pthread -o prog prog.c
\end{verbatim}

% ============================================================================
\section{Platform Support}
% ============================================================================

\begin{table}[h]
\centering
\begin{tabular}{@{}llll@{}}
\toprule
\textbf{Platform} & \textbf{Backtrace} & \textbf{Break} & \textbf{Notes} \\
\midrule
Linux & Yes (execinfo.h) & Yes (\_\_builtin\_trap) & Use \texttt{-rdynamic} \\
macOS & Yes (execinfo.h) & Yes (\_\_builtin\_trap) & Full support \\
Windows (MSVC) & Yes (dbghelp) & Yes (\_\_debugbreak) & Link dbghelp.lib \\
Windows (MinGW) & Fallback & Yes (\_\_builtin\_trap) & Limited backtrace \\
\bottomrule
\end{tabular}
\caption{Platform support matrix}
\end{table}

% ============================================================================
\section{Complete Example}
% ============================================================================

\begin{lstlisting}
/*
 * example.c - Complete debug.h usage example
 */

#define DEBUG_LEVEL 5
#define DEBUG_TIMESTAMPS
#define DEBUG_COLORS
#include "debug.h"
#include <string.h>

typedef struct {
    int id;
    char name[32];
    double score;
} Student;

Student* create_student(int id, const char* name, double score) {
    DEBUG_ENTER("(id=%d, name=%s, score=%.2f)", id, name, score);
    
    Student* s = (Student*)DEBUG_MALLOC(sizeof(Student));
    DEBUG_ASSERT(s != NULL, "Failed to allocate Student");
    
    s->id = id;
    strncpy(s->name, name, sizeof(s->name) - 1);
    s->name[sizeof(s->name) - 1] = '\0';
    s->score = score;
    
    DEBUG_PTR(s, "Student struct");
    DEBUG_VAR(s->id);
    DEBUG_VAR(s->score);
    
    DEBUG_EXIT(s);
}

void process_students(Student** students, int count) {
    DEBUG_ENTER("(count=%d)", count);
    DEBUG_CHECK(count > 0, "No students to process");
    
    double total = 0.0;
    
    DEBUG_LOOP_START("Calculating average");
    for (int i = 0; i < count; i++) {
        DEBUG_LOOP_ITER(i, count);
        total += students[i]->score;
    }
    DEBUG_LOOP_END();
    
    double average = total / count;
    DEBUG_VAR(average);
    DEBUG_INFO("Average score: %.2f", average);
}

int main(void) {
    Student* s1 = create_student(1, "Alice", 95.5);
    Student* s2 = create_student(2, "Bob", 87.3);
    
    Student* students[] = {s1, s2};
    process_students(students, 2);
    
    DEBUG_FREE(s1);
    DEBUG_FREE(s2);
    
    return 0;
}
\end{lstlisting}

% ============================================================================
\section{API Reference Summary}
% ============================================================================

\begin{table}[h]
\centering
\small
\begin{tabular}{@{}ll@{}}
\toprule
\textbf{Macro} & \textbf{Description} \\
\midrule
\multicolumn{2}{l}{\textit{Logging}} \\
\texttt{DEBUG\_TRACE/DEBUG/INFO/WARN/ERROR} & Log at specified level \\
\texttt{DEBUG\_LOG(level, fmt, ...)} & Log at explicit level \\
\midrule
\multicolumn{2}{l}{\textit{Function Tracing}} \\
\texttt{DEBUG\_ENTER(fmt, ...)} & Log function entry \\
\texttt{DEBUG\_EXIT(ret)} & Log exit with return \\
\midrule
\multicolumn{2}{l}{\textit{Inspection}} \\
\texttt{DEBUG\_VAR(var)} & Print variable value \\
\texttt{DEBUG\_PTR(ptr, desc)} & Print pointer info \\
\midrule
\multicolumn{2}{l}{\textit{Memory}} \\
\texttt{DEBUG\_MALLOC(size)} & Logged malloc \\
\texttt{DEBUG\_FREE(ptr)} & Logged free \\
\midrule
\multicolumn{2}{l}{\textit{Assertions}} \\
\texttt{DEBUG\_ASSERT(expr, ...)} & Hard assert (aborts) \\
\texttt{DEBUG\_CHECK(expr, ...)} & Soft check (continues) \\
\texttt{DEBUG\_BACKTRACE()} & Print stack trace \\
\midrule
\multicolumn{2}{l}{\textit{Loops}} \\
\texttt{DEBUG\_LOOP\_START/ITER/END} & Loop tracing \\
\midrule
\multicolumn{2}{l}{\textit{Output}} \\
\texttt{DEBUG\_SET\_OUTPUT(file)} & Redirect to file \\
\texttt{DEBUG\_RESET\_OUTPUT()} & Reset to stderr \\
\texttt{DEBUG\_SET\_HANDLER(fn)} & Custom handler \\
\midrule
\multicolumn{2}{l}{\textit{GDB}} \\
\texttt{DEBUG\_BREAK()} & Insert breakpoint \\
\texttt{DEBUG\_IF(expr)} & Conditional block \\
\texttt{DEBUG\_ABORT(msg, ...)} & Abort with message \\
\bottomrule
\end{tabular}
\caption{Complete API reference}
\end{table}

% ============================================================================
\section{License}
% ============================================================================

MIT License

Copyright (c) 2025 EdgeOfAssembly

Permission is hereby granted, free of charge, to any person obtaining a copy
of this software and associated documentation files (the ``Software''), to deal
in the Software without restriction, including without limitation the rights
to use, copy, modify, merge, publish, distribute, sublicense, and/or sell
copies of the Software, and to permit persons to whom the Software is
furnished to do so, subject to the following conditions:

The above copyright notice and this permission notice shall be included in all
copies or substantial portions of the Software.

THE SOFTWARE IS PROVIDED ``AS IS'', WITHOUT WARRANTY OF ANY KIND, EXPRESS OR
IMPLIED, INCLUDING BUT NOT LIMITED TO THE WARRANTIES OF MERCHANTABILITY,
FITNESS FOR A PARTICULAR PURPOSE AND NONINFRINGEMENT. IN NO EVENT SHALL THE
AUTHORS OR COPYRIGHT HOLDERS BE LIABLE FOR ANY CLAIM, DAMAGES OR OTHER
LIABILITY, WHETHER IN AN ACTION OF CONTRACT, TORT OR OTHERWISE, ARISING FROM,
OUT OF OR IN CONNECTION WITH THE SOFTWARE OR THE USE OR OTHER DEALINGS IN THE
SOFTWARE.

\end{document}
